\[
(\mathbf{X}-\boldsymbol{\mu})'\boldsymbol{\Sigma}^{-1}(\mathbf{X}-\boldsymbol{\mu})
\]

是从 $\mathbf{X}$ 到总体均值向量 $\boldsymbol{\mu}$ 的统计距离平方。如果一个分量有比其他分量大得多的方差，它将对距离平方影响较小。再者，两个高度相关的随机变量比两个几乎不相关的变量影响更小。本质上，利用协方差矩阵的逆，(1)使所有的变量标准化，(2)消除相关的影响。从结论 4.7 的证明，

\[
(\mathbf{X}-\boldsymbol{\mu})'\boldsymbol{\Sigma}^{-1}(\mathbf{X}-\boldsymbol{\mu})=Z_1^2+Z_2^2+\cdots+Z_p^2
\]

依据 $\boldsymbol{\Sigma}^{-1/2}$（参见\hyperref[eq:2-22]{式(2-22)}），$\mathbf{Z}=\boldsymbol{\Sigma}^{-1/2}(\mathbf{X}-\boldsymbol{\mu})$ 有 $N_p(\mathbf{0},I_p)$ 分布，且

\begin{align*}
(\mathbf{X}-\boldsymbol{\mu})'\boldsymbol{\Sigma}^{-1}(\mathbf{X}-\boldsymbol{\mu})&=(\mathbf{X}-\boldsymbol{\mu})'\boldsymbol{\Sigma}^{-1/2}\boldsymbol{\Sigma}^{-1/2}(\mathbf{X}-\boldsymbol{\mu})\\
&=\mathbf{Z}'\mathbf{Z}=Z_1^2+Z_2^2+\cdots+Z_p^2
\end{align*}

统计距离平方似乎可以计算如下：首先，把随机向量 $\mathbf{X}$ 变换为 $p$ 个独立的标准正态随机变量，然后应用通常的距离平方，即变量的平方和。

下面，考虑向量随机变量的线性组合

\phantomsection\label{eq:4-10}
\begin{equation}
c_1\mathbf{X}_1+c_2\mathbf{X}_2+\cdots+c_n\mathbf{X}_n=[\mathbf{X}_1:\mathbf{X}_2:\cdots:\mathbf{X}_n]_{(p\times n)}\mathbf{c}_{(n\times1)}
\tag{4-10}
\end{equation}

这个线性组合不同于以前考虑过的线性组合，原因在于：它定义的一个 $p\times1$ 向量随机变量是向量的一个线性组合。以前，我们讨论过单一的随机变量能写成其他一元随机变量的线性组合。

\mvthm{结论 4.8}{设 $\mathbf{X}_1,\mathbf{X}_2,\cdots,\mathbf{X}_n$ 是相互独立的，且 $\mathbf{X}_j$ 的分布为 $N_p(\boldsymbol{\mu}_j,\boldsymbol{\Sigma})$（注意，每个 $\mathbf{X}_j$ 有相同的协方差矩阵 $\boldsymbol{\Sigma}$）。则}

\[
V_1=c_1\mathbf{X}_1+c_2\mathbf{X}_2+\cdots+c_n\mathbf{X}_n
\]

服从 $N_p\left(\sum\limits_{j=1}^{n}c_j\boldsymbol{\mu}_j,\left(\sum\limits_{j=1}^{n}c_j^2\right)\boldsymbol{\Sigma}\right)$ 分布，而且 $V_1$ 和 $V_2=b_1\mathbf{X}_1+b_2\mathbf{X}_2+\cdots+b_n\mathbf{X}_n$ 是协方差矩阵为

\[
\begin{bmatrix}
\left(\sum\limits_{j=1}^{n}c_j^2\right)\boldsymbol{\Sigma} & (b'c)\boldsymbol{\Sigma}\\[6pt]
(b'c)\boldsymbol{\Sigma} & \left(\sum\limits_{j=1}^{n}b_j^2\right)\boldsymbol{\Sigma}
\end{bmatrix}
\]

的联合多元正态。所以，若 $b'c=\sum\limits_{j=1}^{n}c_jb_j=0$，则 $V_1$ 与 $V_2$ 独立。

\noindent\textbf{证明}\quad 由结论 4.5(c)，$np$ 个分量向量

\[
[X_{11},\cdots,X_{1p},X_{12},\cdots,X_{2p},\cdots,X_{np}]=[\mathbf{X}_1',\mathbf{X}_2',\cdots,\mathbf{X}_n']_{(1\times np)}=\mathbf{X}_{(np\times1)}'
\]

为多元正态。特别的，$\mathbf{X}$ 服从 $N_{np}(\boldsymbol{\mu},\boldsymbol{\Sigma}_x)$ 分布，其中

\[
\boldsymbol{\mu}_{(np\times1)}=\begin{bmatrix}\boldsymbol{\mu}_1\\\boldsymbol{\mu}_2\\\vdots\\\boldsymbol{\mu}_n\end{bmatrix},\qquad
\boldsymbol{\Sigma}_{x\ (np\times np)}=\begin{bmatrix}
\boldsymbol{\Sigma} & \mathbf{0} & \cdots & \mathbf{0}\\
\mathbf{0} & \boldsymbol{\Sigma} & \cdots & \mathbf{0}\\
\vdots & \vdots & \ddots & \vdots\\
\mathbf{0} & \mathbf{0} & \cdots & \boldsymbol{\Sigma}
\end{bmatrix}
\]

选择

\[
A_{(2p\times np)}=\begin{bmatrix}c_1I & c_2I & \cdots & c_nI\\ b_1I & b_2I & \cdots & b_nI\end{bmatrix}
\]

其中 $I$ 是 $p\times p$ 单位矩阵，得到
